\documentclass[11pt,a4paper]{article}
\usepackage[T1]{fontenc}
\usepackage[utf8]{inputenc}
\usepackage{amsmath,amssymb,amsthm}
\usepackage[margin=1in]{geometry}
\usepackage[colorlinks=true,linkcolor=blue,urlcolor=blue]{hyperref}
\theoremstyle{plain}
\newtheorem{theorem}{Theorem}
\newtheorem{lemma}[theorem]{Lemma}
\newtheorem{proposition}[theorem]{Proposition}
\newtheorem{corollary}[theorem]{Corollary}
\theoremstyle{definition}
\newtheorem{definition}[theorem]{Definition}
\newtheorem{remark}[theorem]{Remark}
\title{A proof of the conjectured recurrence for OEIS A258547}
\author{Adrian Perez Fontelles\\ \small Independent researcher}
\date{30 August 2026}
\begin{document}
\maketitle

\begin{abstract}
OEIS A258547 carries a conjectured linear recurrence with polynomial coefficients, of
order 3. The entry posts no generating function, so the usual route through a
functional equation is unavailable. It does post a closed form for $a(n)$, and that is
enough: the closed form is a sum of finitely many hypergeometric terms, and for such a
term every shift quotient is a rational function of $n$. The conjectured recurrence
therefore reduces to two rational-function identities, one for each similarity
class of terms, and each is decided exactly by cancellation. No expansion is truncated and
nothing is guessed; the procedure returns a proof or a refutation. Here it returns a proof,
valid for every $n$ outside an explicitly computed finite set.
\end{abstract}

\noindent\small 2020 Mathematics Subject Classification. 33F10, 11B37, 05A10.\normalsize

\section{The sequence and the conjecture}

OEIS A258547 is ``Number of (n+1)X(1+1) 0..1 arrays with every 2X2 subblock ne-sw antidiagonal difference nondecreasing horizontally and nw+se diagonal sum nondecreasing vertically''. It has offset $1$ and begins
\[
a(1),\dots,a(8)\;=\;16, 44, 104, 228, 480, 988, 2008, 4052,\ \dots
\]
The entry states, not as a conjecture,
\begin{quote}\small
a(n) = 16*2\textasciicircum{}n - 4*n - 12.
\end{quote}
that is,
\begin{equation}\label{eq:cf}
a(n)\;=\;16 \cdot 2^{n} - 4 n - 12 .
\end{equation}
This is the only input the proof takes from the entry, and it was checked against the
entry's own DATA before use.

The entry separately carries the following comment:
\begin{quote}\small
Empirical: a(n) = 4*a(n-1) -5*a(n-2) +2*a(n-3)
\end{quote}
As of the ``Last modified'' line on the live entry (2025-07-23, revision 8) the
statement is still recorded as a conjecture, and no proof appears among the entry's links,
formulas or programs.

Written out, the conjecture is
\begin{equation}\label{eq:conj}
1\,a(n) - 4\,a(n-1) + 5\,a(n-2) - 2\,a(n-3)
\end{equation}

\section{Hypergeometric terms}

\begin{definition}
A nonzero expression $c(n)$ is a \emph{hypergeometric term} in $n$ if the quotient
$c(n+1)/c(n)$ is a rational function of $n$. Two hypergeometric terms $c$ and $d$ are
\emph{similar}, written $c\sim d$, if $c(n)/d(n)$ is a rational function of $n$.
\end{definition}

Similarity is an equivalence relation on hypergeometric terms. Products of factorials,
binomial coefficients with arguments linear in $n$, exponentials $z^{n}$ and polynomials
are hypergeometric, and so is any product of these; a sum of them need not be.

\begin{lemma}\label{lem:shift}
If $c$ is a hypergeometric term then for every integer $i\ge 0$ the shift quotient
$c(n-i)/c(n)$ is a rational function of $n$, computable from $c(n+1)/c(n)$ by a finite
number of substitutions and multiplications.
\end{lemma}

\begin{proof}
Write $\rho(n)=c(n+1)/c(n)$, a rational function by hypothesis. Then
$c(n-1)/c(n)=1/\rho(n-1)$, and inductively
$c(n-i)/c(n)=\prod_{j=1}^{i}\rho(n-j)^{-1}$, a finite product of rational functions.
\end{proof}

The fact that makes the test complete rather than merely sufficient is the linear
independence of dissimilar terms. It is classical; see Petkov\v{s}ek, Wilf and Zeilberger
\cite{AB}, Chapter~5.

\begin{lemma}\label{lem:indep}
Let $c_{1},\dots,c_{m}$ be hypergeometric terms, pairwise dissimilar. If
$\sum_{j=1}^{m}u_{j}(n)c_{j}(n)=0$ for rational functions $u_{j}$ and all large $n$, then
every $u_{j}$ is the zero rational function.
\end{lemma}

\section{The recurrence as a finite set of rational identities}

\begin{proposition}\label{prop:main}
Let $a(n)=\sum_{j=1}^{m}c_{j}(n)$ with each $c_{j}$ hypergeometric, and let
$p_{0},\dots,p_{r}$ be polynomials. Group the $c_{j}$ into similarity classes and write
each class as $g_{\ell}(n)\,f_{\ell,s}(n)$ for a representative $g_{\ell}$ and rational
functions $f_{\ell,s}$. Put
\[
S_{\ell}(n)\;=\;\sum_{i=0}^{r}\;\sum_{s}\;p_{i}(n)\,f_{\ell,s}(n-i)\,
\frac{g_{\ell}(n-i)}{g_{\ell}(n)} ,
\]
a rational function of $n$ by Lemma~\ref{lem:shift}. Then
\[
\sum_{i=0}^{r}p_{i}(n)\,a(n-i)\;=\;\sum_{\ell}g_{\ell}(n)\,S_{\ell}(n),
\]
and the recurrence $\sum_{i}p_{i}(n)a(n-i)=0$ holds for all large $n$ if and only if
$S_{\ell}=0$ for every $\ell$.
\end{proposition}

\begin{proof}
The displayed identity is the definition of $S_{\ell}$ rearranged: each $a(n-i)$ is the
sum of its terms, each term is its class representative times a rational function, and the
representatives are pulled out. The representatives $g_{\ell}$ are pairwise dissimilar by
construction, so Lemma~\ref{lem:indep} applies to the right-hand side: it vanishes for all
large $n$ exactly when every $S_{\ell}$ vanishes identically. Sufficiency is immediate.
\end{proof}

Proposition~\ref{prop:main} turns the conjecture into a finite computation in
$\mathbb{Q}(n)$, with no truncation and no search. The only $n$ it says nothing about are
those where some $S_{\ell}$ has a pole or some $c_{j}(n-i)$ is undefined, and these are
finite in number and listed below.

\section{The computation for A258547}

The closed form \eqref{eq:cf} splits into two similarity classes, with
representativees
\[
-12,\qquad 16 \cdot 2^{n}
\]
Carrying out Proposition~\ref{prop:main} with the coefficients of \eqref{eq:conj} gives
\[
S_{1}(n)\;=\;0,\qquad S_{2}(n)\;=\;0
\]
so every class residual vanishes identically and the conjecture follows.

\begin{theorem}
For all $n\ge 4$,
\[
1\,a(n) - 4\,a(n-1) + 5\,a(n-2) - 2\,a(n-3)
\]
\end{theorem}

\begin{proof}
Immediate from Proposition~\ref{prop:main} and the computation above, the excluded values
being those at which a class residual has a pole or a term of \eqref{eq:cf} is undefined.
\end{proof}

\section{Verification}

Two checks were made, and both are independent of the algebra above.

First, the closed form \eqref{eq:cf} was evaluated at the first $13$ indices and
compared term by term with the entry's DATA; the values agree exactly. A formula that did
not reproduce the entry's own terms would not have been used.

Second, the conjectured recurrence \eqref{eq:conj} was evaluated on the published terms in
exact integer arithmetic at every index where it applies: $25$ instances beginning at
$n=4$, all of them zero. This is not part of the proof, which is symbolic, but it
would have caught a transcription error in the coefficients.

\begin{thebibliography}{9}
\bibitem{AB} M.~Petkov\v{s}ek, H.~S.~Wilf and D.~Zeilberger, \emph{A=B}, A K Peters,
Wellesley MA, 1996.
\bibitem{OEIS} OEIS Foundation Inc., \emph{The On-Line Encyclopedia of Integer Sequences},
entry \href{https://oeis.org/A258547}{A258547}, 2025-07-23.
\bibitem{GKP} R.~L.~Graham, D.~E.~Knuth and O.~Patashnik, \emph{Concrete Mathematics},
2nd ed., Addison-Wesley, 1994.
\end{thebibliography}

\end{document}
